\documentclass[10pt]{article}
\usepackage[T1]{fontenc}
\usepackage{lmodern,amsmath,amssymb,amsthm,microtype,booktabs}
\usepackage[margin=0.84in]{geometry}
\usepackage[colorlinks=true,urlcolor=blue,citecolor=blue]{hyperref}
\newtheorem{theorem}{Theorem}
\newtheorem{lemma}{Lemma}
\newtheorem{corollary}{Corollary}
\newcommand{\N}{\mathbb N}
\title{Recurrence of Greedy Valuation Walks\\and the OEIS A382357 Permutation Conjecture}
\author{Qichao Wang\\{\small Hebei University of Technology}\\[0.3em]
Daoyu Dong\\{\small University of Electronic Science and Technology of China}}
\date{October 26, 2026}
\hypersetup{pdftitle={Recurrence of Greedy Valuation Walks and the OEIS A382357 Permutation Conjecture},pdfauthor={Qichao Wang and Daoyu Dong}}
\begin{document}
\maketitle
\begin{abstract}
We prove that the greedy sequence of distinct positive integers whose
successive $2$-adic valuations differ by one is a permutation of the
positive integers, resolving the conjecture in OEIS A382357. More
generally, for every integer base $b\ge2$, every fixed valuation jump
$d\ge1$, and every finite legal seed, the resulting sequence lists
every integer in the compatible valuation residue class. A last-exit
argument converts hypothetical escape to large valuations into a
contracting recurrence for positive integer edge-crossing counts.
For the binary sequence we also prove the empirical pairing formula
in A382360 and an exact five-term block around each new even power of two.
\end{abstract}

\section{Definitions and the general statement}
For an integer $b\ge2$, define $v_b(n)=\max\{k\ge0:b^k\mid n\}$.
This definition also applies when $b$ is composite. Fix an integer
$d\ge1$ and a finite seed of distinct positive integers
$a_1,\ldots,a_T$ satisfying
\[
 |v_b(a_{i+1})-v_b(a_i)|=d\qquad(1\le i<T).
\]
Thereafter choose $a_{t+1}$ to be the smallest unused positive integer
whose valuation differs from $v_b(a_t)$ by exactly $d$. There is always
an available term, since the higher neighboring valuation class is infinite.
All valuations belong to a single residue class $r\pmod d$, $0\le r<d$.
Let
\[
 \mathcal C_{b,d,r}=\{n\ge1:v_b(n)\equiv r\pmod d\}.
\]

\begin{theorem}[Recurrence and exhaustive coverage]\label{thm:main}
For every $b,d$ and finite legal seed above, the full sequence consisting
of the seed and its greedy continuation lists every member of
$\mathcal C_{b,d,r}$ exactly once.
In particular, with $b=2$, $d=1$, and $a_1=1$, it proves that OEIS
A382357 is a permutation of the positive integers.
\end{theorem}

The OEIS entry \cite{a382357} records the permutation assertion as a
conjecture. It already proves that every power of two appears. The
argument below proves recurrence at every valuation level and therefore
coverage of all integers in each class.

\section{Queues, last exits, and a contraction}
Normalize a valuation $r+dj$ to the level $j\ge0$, so the sequence walks
on the nonnegative integers by steps $\pm1$. Put $\lambda=b^d$ and
\[
 h_b(t)=t+\left\lfloor\frac{t-1}{b-1}\right\rfloor\qquad(t\ge1).
\]
These are, in increasing order, the positive integers not divisible by $b$.
At level $j$ the labels are $b^r\lambda^j h_b(t)$, with $t=1,2,\ldots$.
Each selection uses the smallest unused label at that level. Thus, at
a level absent from the seed, after $c_j$ visits the next label is
\begin{equation}\label{eq:queue}
 Q_j=b^r\lambda^j h_b(c_j+1).
\end{equation}

\begin{lemma}[Propagation of infinite visits]\label{lem:propagate}
If one level is visited infinitely often, then every level is visited
infinitely often, and every compatible integer appears.
\end{lemma}
\begin{proof}
Suppose a level $j$ is visited infinitely often and one neighboring
level has only finitely many visits. The latter has a fixed smallest
unused label $M$ after its last visit. At every sufficiently late visit
to $j$, the next chosen label is at most $M$. Distinctness permits only
finitely many such choices, a contradiction. Infinite visits thus
propagate across every edge of the level path. At any level, a fixed
unused label would bound every subsequent selection from that level,
again contradicting infinitely many distinct selections.
\end{proof}

\begin{lemma}[Last-exit inequality]\label{lem:last}
Suppose every level is visited only finitely many times. Let $K$ be
the greatest level appearing in the seed, and let $U_j$ be the total
number of upward crossings from $j$ to $j+1$. Then every $U_j$ with
$j\ge K$ is a positive integer, and for every $k\ge K+2$,
\begin{equation}\label{eq:contract}
 \lambda^2 U_k\le U_{k-2}+U_{k-1}
                 +\lambda^2-\lambda^2/b-1.
\end{equation}
\end{lemma}
\begin{proof}
Since each finite set of levels receives finitely many visits, the
level of the walk tends to infinity. For $j\ge K$, the walk starts
below or on $j$ and eventually stays above it. Thus the downward
crossing count is $D_j=U_j-1$, and $U_j\ge1$. At a level $j>K$, which
does not contain the initial term, every visit is an arrival; hence its
total visit count is
\begin{equation}\label{eq:flow}
 c_j=U_{j-1}+D_j=U_{j-1}+U_j-1.
\end{equation}

Consider the final departure from level $k\ge K+2$. It must go upward:
a downward departure would require a later crossing of $k$ on the way
to infinity. At this instant level $k-1$ has already had all its
$c_{k-1}$ visits. Level $k+1$ has had at least $U_k-1$ visits, from the
earlier upward crossings of this edge. Both levels are absent from
the seed. The greedy comparison using \eqref{eq:queue} therefore gives
\begin{equation}\label{eq:compare}
 \lambda^2 h_b(U_k)<h_b(c_{k-1}+1).
\end{equation}
The inequality is strict because labels from distinct valuation classes
are different.

Write $\alpha=b/(b-1)$. The explicit formula for $h_b$ gives
\[
 h_b(t)\ge\alpha t-1,\qquad h_b(c+1)\le\alpha c+1.
\]
Consequently \eqref{eq:compare} and \eqref{eq:flow} imply
\begin{align*}
 \lambda^2 U_k
 &<c_{k-1}+\frac{\lambda^2+1}{\alpha}\\
 &=U_{k-2}+U_{k-1}+\lambda^2-\lambda^2/b-1/b.
\end{align*}
The quantities other than $1/b$ are integers, since $\lambda=b^d$.
Rounding down gives \eqref{eq:contract}.
\end{proof}

\begin{proof}[Proof of Theorem~\ref{thm:main}]
If every level has finite visits, Lemma~\ref{lem:last} applies.
For a positive integer $M$, define
\[
 F(M)=\left\lfloor
 \frac{2M+\lambda^2-\lambda^2/b-1}{\lambda^2}\right\rfloor.
\]
We have $F(M)\le M-1$: indeed
\[
 (\lambda^2-2)M-\lambda^2+\lambda^2/b+1
 \ge\lambda^2/b-1>0.
\]
If two consecutive crossing counts are at most $M$, the next two are
at most $F(M)$, by \eqref{eq:contract}. Starting with
$M=\max(U_K,U_{K+1})$, after at most $M$ such pairs the positive
crossing counts are bounded by zero. This contradiction proves that
some level has infinitely many visits. Lemma~\ref{lem:propagate} then
gives every compatible integer, and distinctness is part of the rule.
\end{proof}

For the canonical binary sequence the central inequality reduces to
\[
 4U_k\le U_{k-2}+U_{k-1}+1.
\]
Its incompatibility with an infinite positive integer sequence is the
obstruction to escaping permanently to larger $2$-adic valuations.

\section{Paired powers of two and exact local blocks}
Return to the canonical binary sequence $a(1),a(2),\ldots$.
Let $t_j$ be the index of the unique occurrence of $2^j$. These indices
are OEIS A382360 \cite{a382360}.

\begin{theorem}[Pairing and a five-term block]\label{thm:pairs}
For every $m\ge0$,
\begin{equation}\label{eq:pairs}
 t_{2m+1}=t_{2m}+1.
\end{equation}
More precisely, for $m\ge1$ and $x=4^m$, with $t=t_{2m}$,
\begin{equation}\label{eq:block}
 (a(t-1),a(t),a(t+1),a(t+2),a(t+3))
   =(3x/2,\ x,\ 2x,\ 3x,\ 5x/2).
\end{equation}
The initial block is $(a(1),a(2),a(3))=(1,2,3)$.
\end{theorem}
\begin{proof}
The first visit to level $j$ uses $2^j$, because labels within that
level are used in increasing order. Initially the walk visits levels
$0,1,0$, with labels $1,2,3$.

Suppose the first visits to levels $2m-2$ and $2m-1$ are consecutive,
and are followed by a return to level $2m-2$. The next visit to
$2m-1$ is its second visit, with label $3\cdot2^{2m-1}$. At that
time level $2m-2$ has been visited at least twice. Its next label is
at least $5\cdot2^{2m-2}$, greater than the unused label
$2^{2m}=4\cdot2^{2m-2}$. Thus this second visit immediately enters
the new level $2m$.

Put $x=2^{2m}$. Level $2m-1$ has now been visited exactly twice.
After choosing $x$, the lower candidate is $5x/2$ and the higher one
is $2x$, so $2x$ is chosen. The next candidates are $3x$ and $4x$,
so the walk returns to level $2m$, using $3x$. Its candidates are
then $5x/2$ and $6x$, so $5x/2$ is chosen. This proves
\eqref{eq:block}, supplies the consecutive new levels and immediate
return needed for the next induction step, and establishes
\eqref{eq:pairs}. All visits used in this induction exist by
Theorem~\ref{thm:main}.
\end{proof}

The OEIS formula for A382360 is currently marked empirical; equation
\eqref{eq:pairs} proves it. The five-term identity gives additional
local values without requiring a closed formula for the positions $t_{2m}$.

\section{Reproducible checks}
The proof is independent of finite computation. The supplied verifier
compares a queue implementation with an independent direct scan for
the smallest unused eligible integer, including arbitrary finite seeds
and bases and jumps other than two and one. It checks the displayed
OEIS prefix, twenty recorded inverse positions, the local blocks, and
the crossing inequalities on finite first-arrival paths. Separate
integer audits test the rounding and contraction estimates.

\begin{thebibliography}{9}
\bibitem{a382357} R. Sigrist, OEIS A382357,
\url{https://oeis.org/A382357}, accessed October 26, 2026.
\bibitem{a382360} R. Sigrist, OEIS A382360,
\url{https://oeis.org/A382360}, accessed October 26, 2026.
\end{thebibliography}
\end{document}
